\section{讨论}
\label{sec:discussion}
\subsection{既有问题中的经验诊断}
图例相容性、异质标注和类别空间差异均有既有研究基础\citep{yang2017harmonization,lambert2020mseg,kim2022uniseg}。TACOSS 与 LandSegmenter 也研究可迁移土地覆盖语义\citep{zermatten2025tacoss,liu2026landsegmenter}。本文的贡献是有明确边界的经验审计：把声明的类别映射关联到特定学习器实际可用的掩膜，再在指定控制下报告所考察源混合的恢复与源支持类别损失。本文不主张首次提出类别体系构建或发现负迁移。

精确标签缺失限定了可用监督，但并不排除影像偏移，不能确定失败的唯一原因，也不意味着预训练或异质监督模型无法识别该概念。主要 CLIP 对比和历史原生 U-Net 对比在各自记录的训练制度下呈现相同方向。这一扩展缩小了仅归因于单一 CLIP 校准工具的解释范围，但同一目标上的两个选定模型家族并不足以证明架构层面的普适性。

\subsection{源混合与配对能够支持什么结论}
主要实验同时加入 FLAIR 影像和草本标签。暴露控制与重复裁块控制约束两种替代解释，但不能将该对比变为仅标签干预。已记录的错配配置保留大部分恢复，原生 B1 只进一步增加较小差值。缺少完整主要排程限制了配对与暴露约束的独立核验。辅助影像内容、区域统计以及同区域错排中的弱对齐仍是相容解释。保留恢复比例只是算术描述，并非归因于某组件的因果百分比。已停用的 U-Net 与反向源 shuffle 分支不提供进一步效果证据。

缺失类别获得收益并不等于总体迁移更好，也不等于保留既有能力。适应研究已报告共同与私有类别性能，并考察已知类别负迁移\citep{choe2024bus,choe2025unida,ren2026sats}。本文的损失涉及六个 LoveDA 精确支持类别，或反向案例中的三个 Semantic Drone 精确支持类别。决策边界变化、源外观、采样与损失权重均可能参与其中。实验未证明表征容量是原因，也未验证缓解模块。

\subsection{复现与来源证据边界}
Semantic Drone 使用主要 CLIP 检查点提供第二目标评估，但其草地对其余像素指标不衡量全部支持类别保持。反向源案例改变缺失类别，却复用 OpenEarthMap 目标版本。原生 U-Net 记录改变架构，但保持主要目标总体。这些是共享部分组件的互补检验，而非完整独立的架构与目标交叉复现。B2 证明另一直接监督制度下的可学习性，不能证明外部迁移优越性。

Semantic Drone 三个检查点的点值可以从历史整数总计重构，但两个检查点的逐栅格记录未找回，因此不使用历史多检查点区间。原生 U-Net 标识符集合和逐栅格参考边际与主要总体一致，并不因此验证空间标签相同或原检查点实际执行。源就绪审计内容通过哈希恢复，不恢复历史时间戳或实际执行代码身份。城市与栅格区间以记录检查点为条件，不代表训练种子不确定性，也不证明 Semantic Drone 地理站点独立。地理排除降低已知重叠风险，却不能证明坐标级分离。在上述边界内，部署报告仍可明确呈现所考察缺失类别收益及其源支持类别损失。

\subsection{类别代理与构念效度边界}
固定编码表不会使原生标注意义自动相同。OEM rangeland 包含灌丛与灌木，因此分配到草本输出时只是较宽的植被代理（英文补充材料 S2）。报告的 rangeland 恢复仍是在声明映射下的原生标签得分，不证明仅对纯草本像素的恢复。保留的 OEM 参考标签不提供草本与灌木的细分。因此，本体分辨率和标注范围仍可能与影像、监督共同影响结果。本文保留原始映射、提示词、检查点和整数分数，不合成新的植物学参考，也不通过改变模型选择来处理这一限制。
